% A synthetic manuscript, written to be measured rather than to be read.
%
% Every number, author and result here is invented. The topic -- calibrating
% multiplicative shear bias against image simulations -- is deliberately a
% textbook one with a large published literature, so that nothing in this file
% resembles unpublished work by anyone.
%
% It is written to carry the writing habits the tools measure, at densities a
% referee would notice: an announced enumeration, an introduction with no
% epistemic marking, paragraphs that recite quantities without ranking them,
% adjacent sentences that share no vocabulary, and a few L0 lexical targets.
% Running `tools/ai_ism_lint.py` on it reproduces the README's demo output.
\documentclass{article}

\newcommand{\Nsims}{4200}
\newcommand{\mbias}{0.013}
\newcommand{\cbias}{0.0004}

\begin{document}

\title{Multiplicative Shear Bias in Stacked Image Simulations}
\author{A. Sample}
\maketitle

\begin{abstract}
We measure the multiplicative shear bias of a moments-based shape estimator on
\Nsims{} image simulations. The bias is $m = \mbias$ with an additive term
$c = \cbias$. The calibration holds across signal-to-noise ratio, size, and
blending fraction.
\end{abstract}

\section{Introduction}

Weak lensing surveys measure the coherent distortion of background galaxy
shapes. The shear signal is small and the shape estimator must be calibrated to
a precision far below the statistical error. Image simulations provide that
calibration by constructing a population with known input shear. This work
delivers the calibration for a moments-based estimator.

There are five elements to the approach. First, we build the simulated
population. Second, we match its distribution to the observed catalog. Third,
we run the estimator. Fourth, we fit the bias model. Fifth, we propagate the
residual into the error budget. The pipeline is fully automated.

\section{Method}

The simulation grid spans 12 signal-to-noise bins from 10 to 200, 8 size bins
from 1.2 to 4.6 pixels, and 5 blending fractions from 0.00 to 0.40. Each cell
holds 350 realizations at an input shear of 0.02 applied in 4 position angles.
The total is \Nsims{} images. The point spread function is a Moffat profile
with $\beta = 3.5$ and a full width at half maximum of 0.7 arcsec. The pixel
scale is 0.263 arcsec and the read noise is 4.2 electrons.

The estimator is not a maximum-likelihood fit but a weighted second-moment
measurement --- it measures the quadrupole rather than fitting a profile. The
weight function is a circular Gaussian matched to the object size. It is
important to note that the weight introduces a size-dependent response.

\section{Results}

The measured bias is $m = \mbias \pm 0.002$ and $c = \cbias \pm 0.0002$. In the
lowest signal-to-noise bin $m = 0.041$ and in the highest $m = 0.004$. The
size trend gives $m = 0.028$ below 2 pixels and $m = 0.007$ above 3 pixels. At
a blending fraction of 0.40 the bias reaches 0.052. The residual after
correction is 0.0011 in every cell.

The calibration is stable. The bootstrap scatter over 200 resamples is 0.0008.
Splitting the grid by position angle changes $m$ by 0.0003. A different noise
realization changes it by 0.0005, highlighting the robustness of the result.

\section{Discussion and Conclusion}

The estimator is calibrated to $m = \mbias$. The blending term dominates the
residual budget and future surveys will need deeper simulations. The
photometric redshift distribution was assumed fixed. Detector nonlinearity was
not modelled. The method transfers to any survey with a comparable point
spread function.

\end{document}
